% SPDX-FileCopyrightText: 2016-2026 Lari Natri <lari.natri@iki.fi>
% SPDX-License-Identifier: GPL-3.0-or-later
%
% Repeat-bar dimensions, fragment state/drawing and inline repeat declarations.
% This file is part of the 'ulsbs' package.
%
% Load: ulsbs.sty, after attached notes and before musical navigation marks.
% Requires: init registers, prose content indent, final fonts/style, visibility
% policy, shared boolean/repeat utilities, xparse and expl3. Late verse/music
% owners install lifecycle integration, line metrics and final inline markers.
% State: owns repeat dimensions, nesting/serial counters, continuation map and
% fragment scratch registers. Inherits @ as a letter without a wrapper group;
% retains all original expl3 spans and conditional engine branches.
% ulsbs-spacing.code.tex executes setters/default measurements at the original
% site; the following directlua loads ulsbs-repeats.lua once, allocating the
% attributes and shipout callback before fragment definitions. Public lengths
% may be changed after loading with setlength. ContinuingOverlap expands against
% the current baseline at use time; spacing policy uses its public setters.

\ProvidesFile{ulsbs-repeats.code.tex}[ULSBS repeat bars and inline marks]

% === Line repeat bars ===

% Line repeat bars inside verses. This prints vertical repeat
% bars for the lyric lines within it's effect. Can be nested. Use only within
% verses
%
% Usage example:
%
%   \beginverse
%     lyric line without repeat
%     \beginrep
%       lyric line with one repeat level
%       lyric line with one repeat level (2nd line of the same repeat)
%       \beginrep
%         lyric line with two repeat levels
%       \endrep
%     \endrep
%   \endverse
%
% The verse machinery installed later supplies repeat lifecycle integration.

% Consumer settings: lengths/dimensions are initialized here; override them after
% loading with \setlength. Renew \ulRepeatBarContinuingOverlap after loading;
% unlike the dimensions, it evaluates the current baseline each time it is used.

% Repetition bar width (register provided by init).
\setlength{\ulRepeatBarWidth}{1.4pt}
% The horizontal gap between line repeat bars, and minimum gap between the
% last (deepest level) bar and the start of the lyrics
\newdimen\ulRepeatBarGap\ulRepeatBarGap=2.265pt

% NOTE: Be careful if changing the following settings. Test all situations:
% with chords, lyric-only, without measure bars. Separate bars following each
% must leave enough gap between them, and multiline repeat bars ought to have
% no gaps in between lyric lines.

% Inline lrep/rrep metrics. Verse-left top/bottom geometry uses the policy below.
\newdimen\ulRepeatBarDepth\ulRepeatBarDepth=.25pt
\newdimen\ulRepeatBarHeightAdjustmentWithChords\ulRepeatBarHeightAdjustmentWithChords=-2.5pt
\newdimen\ulRepeatBarHeightAdjustmentNoChords\ulRepeatBarHeightAdjustmentNoChords=.5pt
% At a page/column break, add two short vertical dashes beyond each end of the
% solid bar. These dimensions control dash length and the gaps between the
% solid bar and dashes, and between the two dashes.
\newdimen\ulRepeatBarContinuationDashLength
\ulRepeatBarContinuationDashLength=2pt
\newdimen\ulRepeatBarContinuationDashGap
\ulRepeatBarContinuationDashGap=1.25pt
% How much extra upward reach continuation lines get to bridge interline glue
% when having multiline repeat bars.
\newcommand{\ulRepeatBarContinuingOverlap}{\the\dimexpr .7\baselineskip\relax}

% SPDX-FileCopyrightText: 2016-2026 Lari Natri <lari.natri@iki.fi>
% SPDX-License-Identifier: GPL-3.0-or-later
%
% Bar geometry and verse-spacing setters, validation and initial policy calls.
% This file is part of the 'ulsbs' package.
%
% Load: ulsbs-repeats.code.tex, after repeat dimensions and before repeat state.
% Requires: initialized verse-separation storage, final lyric/verse font hooks,
% xparse and expl3. Music-context synchronization is conditional until the late
% music owner installs it; Lua consumers snapshot expanded numeric macros.
% State: owns validated bar/verse policy macros, scalar/dimension scratch state
% and nominal baseline glue. Inherits @ as a letter without a wrapper group;
% retains the complete original expl3 span and setter-local measurement group.
% Default calls execute eagerly here, in order. Verse setters expand expressions
% once in the configured lyric font and freeze resolved skips; later font changes
% do not rescale them. Public setters are scoped and must follow custom font
% configuration. ulVerseSpacingScale reads the stored scalar policy, not glue.

\ProvidesFile{ulsbs-spacing.code.tex}[ULSBS bar policy and verse spacing]

% Verse-left bar policy. Store expanded numbers, not registers or expressions:
% the line builder snapshots these local values with token.get_macro.
\ExplSyntaxOn
\tl_new:N \l__ulsbs_bar_value_tl
\seq_new:N \l__ulsbs_bar_dimension_seq
\cs_generate_variant:Nn \regex_match:nnTF {nV}
\cs_generate_variant:Nn \regex_extract_once:nnNTF {nV}
\cs_new_protected:Npn \__ulsbs_bar_error:n #1
  {\PackageError{ulsbs}{Invalid~bar~setting:~#1}
    {Use~N=0..3,~L=0..N,~finite~nonnegative~scales,~a~boolean,~or~a~nonnegative~dimension.}}
\prg_new_protected_conditional:Npnn \__ulsbs_bar_scalar:n #1 {TF}
  {
    \tl_set:Nx \l__ulsbs_bar_value_tl {#1}
    \regex_match:nVTF
      {\A\s*[+\-]?(\d+(\.\d*)?|\.\d+)([eE][+\-]?\d+)?\s*\Z}
      \l__ulsbs_bar_value_tl
      % Keep normalized values finite when the Lua consumer reads a double.
      {\fp_compare:nTF {0 <= (\l__ulsbs_bar_value_tl) <= 1.797693134862315e308}
        {\prg_return_true:}{\prg_return_false:}}
      {\prg_return_false:}
  }
\cs_new_protected:Npn \__ulsbs_bar_top:nnnn #1#2#3#4
  {
    \tl_set:Nx \l__ulsbs_bar_value_tl {#2}
    \regex_match:nVTF {\A\s*\+?0*[0-3]\s*\Z} \l__ulsbs_bar_value_tl
      {
        \tl_set:Nx \l__ulsbs_bar_value_tl {#3}
        \regex_match:nVTF {\A\s*\+?0*[0-3]\s*\Z} \l__ulsbs_bar_value_tl
          {\int_compare:nNnTF {#3} > {#2}
            {\__ulsbs_bar_error:n {#1~top~level~must~not~exceed~N}}
            {\__ulsbs_bar_scalar:nTF {#4}
              {
                \cs_set:cpx {ulsbs@#1@level@\int_eval:n{#2}} {\int_eval:n{#3}}
                \cs_set:cpx {ulsbs@#1@scale@\int_eval:n{#2}} {\fp_to_decimal:n{#4}}
              }
              {\__ulsbs_bar_error:n {#1~top~scale~must~be~finite~and~nonnegative}}}}
          {\__ulsbs_bar_error:n {#1~top~level~must~be~an~integer~0..N}}
      }
      {\__ulsbs_bar_error:n {#1~top~N~must~be~an~integer~0..3}}
    \cs_if_exist:NT \__ulsbs_music_context: {\__ulsbs_music_context:}
  }
\NewDocumentCommand \ulMeasureBarTop {mmm}
  {\__ulsbs_bar_top:nnnn {measure}{#1}{#2}{#3}}
\NewDocumentCommand \ulRepeatBarTop {mmm}
  {\__ulsbs_bar_top:nnnn {repeat}{#1}{#2}{#3}}
\NewDocumentCommand \ulUseVerseAnnotationMinimum {m}
  {
    \tl_set:Nx \l__ulsbs_bar_value_tl {\str_lowercase:f {#1}}
    \tl_trim_spaces:N \l__ulsbs_bar_value_tl
    \str_case:VnF \l__ulsbs_bar_value_tl
      {{true}{\cs_set:cpn {ulsbs@bar@verseminimum}{1}}
       {yes}{\cs_set:cpn {ulsbs@bar@verseminimum}{1}}
       {on}{\cs_set:cpn {ulsbs@bar@verseminimum}{1}}
       {1}{\cs_set:cpn {ulsbs@bar@verseminimum}{1}}
       {false}{\cs_set:cpn {ulsbs@bar@verseminimum}{0}}
       {no}{\cs_set:cpn {ulsbs@bar@verseminimum}{0}}
       {off}{\cs_set:cpn {ulsbs@bar@verseminimum}{0}}
       {0}{\cs_set:cpn {ulsbs@bar@verseminimum}{0}}}
      {\__ulsbs_bar_error:n {verse~annotation~minimum~must~be~boolean}}
    \cs_if_exist:NT \__ulsbs_music_context: {\__ulsbs_music_context:}
  }
\NewDocumentCommand \ulRepeatBarBottomScale {m}
  {\__ulsbs_bar_scalar:nTF {#1}
    {\cs_set:cpx {ulsbs@repeat@bottomscale}{\fp_to_decimal:n{#1}}}
    {\__ulsbs_bar_error:n {repeat~bottom~scale~must~be~finite~and~nonnegative}}
    \cs_if_exist:NT \__ulsbs_music_context: {\__ulsbs_music_context:}}
\cs_new_protected:Npn \__ulsbs_bar_dimension:nnn #1#2#3
  {
    \tl_set:Nx \l__ulsbs_bar_value_tl {#1}
    \regex_match:nVTF {\A\s*[+\-]?0+(\.0*)?\s*\Z} \l__ulsbs_bar_value_tl
      {\cs_set:cpn {#2}{0}}
      {
        \regex_extract_once:nVNTF
          {\A\s*([+\-]?(\d+(\.\d*)?|\.\d+))\s*(pt|pc|in|bp|cm|mm|dd|cc|sp|em|ex)\s*\Z}
          \l__ulsbs_bar_value_tl \l__ulsbs_bar_dimension_seq
          {
            % Check the TeX dimension limit before conversion, including in nonstop mode.
            \fp_compare:nTF
              {0 <= (\seq_item:Nn \l__ulsbs_bar_dimension_seq {2}) <=
                1073741823 / (65536 * \dim_to_fp:n {1\seq_item:Nn \l__ulsbs_bar_dimension_seq {5}})}
              {\cs_set:cpx {#2}{\dim_to_decimal_in_sp:n{#1}}}
              {\__ulsbs_bar_error:n {#3~minimum~gap~must~be~nonnegative~and~within~TeX's~dimension~limit}}
          }
          {\__ulsbs_bar_error:n {#3~minimum~gap~requires~a~number~and~dimension~unit}}
      }
    \cs_if_exist:NT \__ulsbs_music_context: {\__ulsbs_music_context:}
  }
\NewDocumentCommand \ulRepeatBarMinimumGap {m}
  {\__ulsbs_bar_dimension:nnn {#1}{ulsbs@repeat@minimumgap}{repeat}}

% Resolve scalar glue once, at the call site, against the configured lyric font.
% Call these setters after custom font configuration; later font changes do not
% rescale the skips, and lyric/source material is never evaluated here. Argument
% expressions expand once inside that font context, before validation or writes.
\skip_new:N \l__ulsbs_verse_nominal_skip
\cs_new_protected:Npn \__ulsbs_verse_error:n #1
  {\PackageError{ulsbs}{Invalid~verse~setting:~#1}
    {Use~three~finite~nonnegative~scales~whose~resolved~dimensions~fit~TeX's~dimension~limit.}}
\cs_new_protected:Npn \__ulsbs_verse_spacing:nnnn #1#2#3#4
  {
    \group_begin:
      \ulLyricFontStyle\ulVerseFont\relax
      \use:x {\group_end:
        \__ulsbs_verse_spacing_apply:nnnnn {#1}{#2}{#3}{#4}
          {\dim_eval:n {\tex_fontcharht:D \tex_font:D `A}}}
  }
\cs_new_protected:Npn \__ulsbs_verse_spacing_apply:nnnnn #1#2#3#4#5
  {
    \__ulsbs_bar_scalar:nTF {#2}
      {\__ulsbs_bar_scalar:nTF {#3}
        {\__ulsbs_bar_scalar:nTF {#4}
          {
            \fp_compare:nTF
              {max(#2,#3,#4) * \dim_to_fp:n {#5}
                <= 1073741823 / 65536}
              {
                \cs_set:cpx {ulsbs@verse@#1@natural}{\fp_to_decimal:n {#2}}
                \cs_set:cpx {ulsbs@verse@#1@stretch}{\fp_to_decimal:n {#3}}
                \cs_set:cpx {ulsbs@verse@#1@shrink}{\fp_to_decimal:n {#4}}
                \skip_set:cn {\str_if_eq:nnTF {#1}{baseline}
                    {l__ulsbs_verse_nominal_skip}{ulsbs@verse@separation}}
                  {\fp_to_dim:n {#2 * \dim_to_fp:n {#5}}
                    plus \fp_to_dim:n {#3 * \dim_to_fp:n {#5}}
                    minus \fp_to_dim:n {#4 * \dim_to_fp:n {#5}}}
              }
              {\__ulsbs_verse_error:n {#1~glue~exceeds~TeX's~dimension~limit}}
          }
          {\__ulsbs_verse_error:n {#1~shrink~must~be~finite~and~nonnegative}}}
        {\__ulsbs_verse_error:n {#1~stretch~must~be~finite~and~nonnegative}}}
      {\__ulsbs_verse_error:n {#1~natural~must~be~finite~and~nonnegative}}
    \cs_if_exist:NT \__ulsbs_music_context: {\__ulsbs_music_context:}
  }
\NewDocumentCommand \ulVerseBaseline {mmm}
  {\__ulsbs_verse_spacing:nnnn {baseline}{#1}{#2}{#3}}
\NewDocumentCommand \ulVerseSeparation {mmm}
  {\__ulsbs_verse_spacing:nnnn {separation}{#1}{#2}{#3}}
\NewDocumentCommand \ulVerseMinimumGap {m}
  {\__ulsbs_bar_dimension:nnn {#1}{ulsbs@verse@minimumgap}{verse}}
\msg_new:nnn {ulsbs}{verse-spacing-key}
  {Invalid~verse~spacing~key~'#1'.~Use~baseline~or~separation,~then~natural,~stretch~or~shrink.}
\cs_new:Npn \__ulsbs_verse_spacing_scale:nn #1#2
  {\str_case:nnF {#2}
    {{natural}{\use:c {ulsbs@verse@#1@natural}}
     {stretch}{\use:c {ulsbs@verse@#1@stretch}}
     {shrink}{\use:c {ulsbs@verse@#1@shrink}}}
    {\msg_expandable_error:nnn {ulsbs}{verse-spacing-key}{#1/#2}}}
\NewExpandableDocumentCommand \ulVerseSpacingScale {mm}
  {\str_case:nnF {#1}
    {{baseline}{\__ulsbs_verse_spacing_scale:nn {baseline}{#2}}
     {separation}{\__ulsbs_verse_spacing_scale:nn {separation}{#2}}}
    {\msg_expandable_error:nnn {ulsbs}{verse-spacing-key}{#1/#2}}}
\ulMeasureBarTop{0}{0}{1.3}
\ulMeasureBarTop{1}{1}{1.3}
\ulMeasureBarTop{2}{2}{.85}
\ulMeasureBarTop{3}{2}{.85}
\ulRepeatBarTop{0}{0}{1.3}
\ulRepeatBarTop{1}{1}{1.3}
\ulRepeatBarTop{2}{2}{.85}
\ulRepeatBarTop{3}{2}{.85}
\ulUseVerseAnnotationMinimum{true}
\ulRepeatBarBottomScale{.025}
\ulRepeatBarMinimumGap{3pt}
\ulVerseBaseline{1.618034}{0}{0}
\ulVerseSeparation{1.618034}{0}{.4}
\ulVerseMinimumGap{1.25pt}
\ExplSyntaxOff


% State: key->boolean map for
%        <repeat nesting level> -> <repeat started on previous line>
\ExplSyntaxOn
\prop_new:N \g__ulsbs_repeat_cont_prop
\ExplSyntaxOff

% State: Current line repeat bar's nesting depth.
% Updated by \beginrep and \endrep.
\newcount\ulsbs@rep@level\ulsbs@rep@level=0\relax%
% Give every repeat occurrence its own id. On LuaTeX the id is attached to
% its rule fragments, which are coalesced into one rule at shipout.
\newcount\ulsbs@rep@serial\ulsbs@rep@serial=0\relax%

% LuaTeX can alter the final node list after layout splits it into physical
% pages and columns. Group tagged fragments by repeat id and absolute x
% position, retain the uppermost rule, and extend it through the lowest one.
% A repeat crossing a page or column boundary is therefore represented by one
% rule in each resulting segment. Other engines keep the fragment fallback.
\ifdefined\directlua
  \directlua{dofile(kpse.find_file('ulsbs-repeats.lua'))}
  \def\ulsbs@rep@tagrule#1#2#3#4{%
    \begingroup%
      \csname __ulsbs_music_context:\endcsname%
      \attribute\ulsbs@rep@attribute=\csname ulsbs@rep@id@#1\endcsname\relax%
      \attribute\ulsbs@rep@heightattribute=\number#3\relax%
      \ifnum#2=1\relax%
        \attribute\ulsbs@rep@startattribute=1\relax%
      \fi%
      #4%
    \endgroup%
  }
  \def\ulsbs@rep@endmarker{%
    \begingroup%
      \attribute\ulsbs@rep@attribute=\csname ulsbs@rep@id@\the\ulsbs@rep@level\endcsname\relax%
      \attribute\ulsbs@rep@endattribute=1\relax%
      \directlua{%
        local marker = node.new("whatsit", "user_defined")
        marker.type = 100
        marker.value = 0
        node.write(marker)
      }%
    \endgroup%
  }
\else
  % Inactive fallback branch: the public engine guard rejects non-Lua engines
  % before this point. Its per-line fragment drawing is not a supported backend;
  % the Lua branch above owns coalescing and page/column continuation dashes.
  \PackageInfo{ulsbs}{Non-Lua engine: repeat bars use per-line fragments;^^J%
    single-object bars and break continuation dashes require LuaTeX}%
  \def\ulsbs@rep@tagrule#1#2#3#4{#4}%
  \let\ulsbs@rep@endmarker\relax
\fi

% Scratch registers:
\newlength\ulsbs@tmp@rep@allbarswidth
\newcount\ulsbs@tmp@rep@prevlinestoreach
\newdimen\ulsbs@tmp@rep@barhtbase
\newdimen\ulsbs@tmp@rep@normalbarhtbase

% Draw one line repetition bar (fragment) for current lyric line. Called
% repetitiously from \ulsbs@rep@drawbars.
\newcommand{\ulsbs@rep@drawbars@iter}[1]{%
  \rlap{%
    \ulIfBool{g__ulsbs_repeat_cont_prop}{#1}{% for the repeat started on earlier lyric line
      % If current bar was started on earlier lyric line (par), increase the
      % current line's repeat bar height to reach through the gap
      \advance\ulsbs@tmp@rep@barhtbase by \ulRepeatBarContinuingOverlap\relax
      % If previous lyric line (par) was (long and) wrapped to more than one
      % line, make the current line's repeat bar higher to reach through
      % them all.
      \ifnum\ulsbs@prevparlines>1\relax%
        \ulsbs@tmp@rep@prevlinestoreach=\ulsbs@prevparlines%
        \advance\ulsbs@tmp@rep@prevlinestoreach by -1\relax%
        \advance\ulsbs@tmp@rep@barhtbase by
         \dimexpr \baselineskip * \numexpr\ulsbs@tmp@rep@prevlinestoreach\relax\relax%
      \fi%
      % Draw the repeat bar with smash to take no vertical space
      \ulsbs@rep@tagrule{#1}{0}{\ulsbs@tmp@rep@normalbarhtbase}{\smash{%
        \vrule\@width\ulRepeatBarWidth%
              \@height\ulsbs@tmp@rep@barhtbase%
              \@depth\ulRepeatBarDepth%
      }}%
    }{%
      % We are on the first line of the repeat: smash or no?
      \ulsbs@rep@tagrule{#1}{1}{\ulsbs@tmp@rep@normalbarhtbase}{\smash{%
        \vrule\@width\ulRepeatBarWidth%
              \@height\ulsbs@tmp@rep@barhtbase%
              \@depth\ulRepeatBarDepth%
      }}%
    }%
    \ulBoolSetTrue{g__ulsbs_repeat_cont_prop}{#1}% The bar at this nesting level has been started
  }%
  \kern\ulRepeatBarWidth%
  \kern\ulRepeatBarGap%
  \addtolength{\ulsbs@tmp@rep@allbarswidth}{\ulRepeatBarWidth}%
  \addtolength{\ulsbs@tmp@rep@allbarswidth}{\ulRepeatBarGap}%
}

% Legacy inline-sign height only; verse fragments use the music owner's policy.
\def\ulsbs@rep@setbarhtbase{%
  \ifulShowChords%
    \global\ulsbs@tmp@rep@barhtbase=\baselineskip%
    \global\advance\ulsbs@tmp@rep@barhtbase \ulRepeatBarHeightAdjustmentWithChords\relax%
  \else%
    \global\settoheight{\ulsbs@tmp@rep@barhtbase}{\vphantom{A}}%
    \global\advance\ulsbs@tmp@rep@barhtbase \ulRepeatBarHeightAdjustmentNoChords\relax%
  \fi%
}

% Draw all repetition bars (fragments) for the current lyric line.
\def\ulsbs@rep@drawbars{%
  {% block to contain color
    \color{ulRepetitionSignColor}%
    \ulsbs@rep@setlinebase%
    \ulsbs@tmp@rep@normalbarhtbase=\ulsbs@tmp@rep@barhtbase%
    \setlength{\ulsbs@tmp@rep@allbarswidth}{0pt}%
    % Call \ulsbs@rep@drawbars@iter for each repeat level
    \ulRepeatTimes{\ulsbs@rep@level}{\ulsbs@rep@drawbars@iter}%
    \ifdim\ulsbs@tmp@rep@allbarswidth<\ulContentIndentWidth%
      \kern\dimexpr\ulContentIndentWidth-\ulsbs@tmp@rep@allbarswidth\relax%
    \fi%
  }%
}

% === Inline repeat marks ===

% The music module owns \lrep/\rrep after the stanza/chord machinery loads.
% Their upper dots follow effective chord policy, never mutable book identity.
% Colorize \rep and use our default font's Unicode × (Noto supports it)
\DeclareDocumentCommand\rep{m}{%
  {\color{ulRepetitionSignColor}(×#1)}%
}

% Use this instead of \rep to signify number of repeats on the first line
% of a line repeat. The parameter is the number of repetitions, as in \rep
% macro.
% Example:
%   \beginverse
%     \beginrep\prep{3}
%       lyric line
%     \endrep
%   \endverse
\newcommand{\prep}[1]{%
  % The music owner replaces this fallback with a retained line-aware marker.
  \raisebox{1ex}{%
    {\color{ulRepetitionSignColor}#1×:}%
  }%
}
